\chapter{Hosting a language of your own}

Chapter~5's front-ends are not privileged. The API they target is public, and
a new language is an ordinary Pop-11 program. This chapter is the recipe;
\code{HOSTING-LANGUAGES.md} in the tree is the longer version, and
\code{ref vmcode} is the specification.

\section{Plant VM code, don't emit machine code}

You write a \textbf{reader}, a \textbf{parser}, and then calls to about forty
\textbf{VM-builder} procedures that plant abstract instructions. The
back-end turns those into native code for whatever CPU you are on.

\begin{center}
\footnotesize
\begin{tabular}{ll}
\toprule
\textbf{To plant\ldots} & \textbf{call}\\
\midrule
a literal constant & \code{sysPUSHQ(item)}\\
the value of a variable & \code{sysPUSH(word)}\\
a store into a variable & \code{sysPOP(word)}\\
a call of a named procedure & \code{sysCALL(word)}\\
a call of a literal procedure & \code{sysCALLQ(item)}\\
start / end of a procedure & \code{sysPROCEDURE(props, nargs)} \ldots \code{sysENDPROCEDURE()}\\
locals, dynamic locals & \code{sysLVARS}, \code{sysLOCAL}, \code{sysNEW\_LVAR}\\
control flow & \code{sysLABEL}, \code{sysGOTO}, \code{sysIFSO}, \code{sysIFNOT}, \code{sysGO\_ON}\\
a permanent assignment & \code{sysPASSIGN(item, token)}\\
a syntax word / macro & \code{sysSYNTAX}\\
\bottomrule
\end{tabular}
\end{center}

The whole model fits in six lines. Here we build, at run time, the equivalent
of \code{define add1(x); x fi\_+ 1 enddefine}:

\begin{pop}
define build_add1() -> p;
    sysPROCEDURE(false, 1);      ;;; open a 1-argument procedure
        sysPUSHQ(1);             ;;;   push literal 1    -> stack: x 1
        sysCALL("fi_+");         ;;;   fast integer add  -> stack: x+1
    sysENDPROCEDURE() -> p;      ;;; p is now a real compiled procedure
enddefine;

vars add1 = build_add1();
isprocedure(add1) =>          ;;; ** <true>
pdnargs(add1)     =>          ;;; ** 1
add1(41)          =>          ;;; ** 42
\end{pop}

\begin{callout}{Execute level}
Planting must happen \emph{while something is running}, not at the top level
of a file being compiled --- at the top level you get
\code{sysEXECUTE: NOT AT EXECUTE LEVEL}. Wrap the plant sequence in a
procedure and call it, as above.
\end{callout}

That is the entire trick. A compiler for your language is a tree walk that
emits the right \code{sys\ldots} calls. What you get for free, without writing
a line of assembly: native code generation on every port, the garbage
collector, the open stack, closures, recursion, the incremental compiler, the
external interface, the REPL, and saved images.

For the reader, you can reuse Pop-11's itemiser (\code{incharitem} over
\code{proglist}, documented in \code{ref popcompile}) with your own character
table --- for a whitespace-delimited language like Forth it needs almost no
configuration --- or write your own, as Prolog does with its
operator-precedence reader.

\section{Three tiers of effort}

\begin{center}
\footnotesize
\begin{tabular}{p{3.1cm}p{7.4cm}l}
\toprule
\textbf{Tier} & \textbf{What you get} & \textbf{Effort}\\
\midrule
1. Embedded interpreter & A plain Pop-11 program: a dictionary and an eval loop. No VM compilation. & days\\
2. VM-compiling front-end & Parse, then plant --- genuinely native definitions, closures, recursion, full interop. \emph{The Poplog way.} & 1--3 weeks\\
3. Full subsystem & Tier 2 plus registration: its own REPL, prompt, banner, file extension and saved image. & +1--3 days\\
\bottomrule
\end{tabular}
\end{center}

The tiers stack. Build tier 1 as a throwaway to pin the semantics with no
compiler in the way, replace its core with planting, then wrap the subsystem
around it.

\section{Tier 3: registering a subsystem}

A subsystem is a record in \code{sys\_subsystem\_table} with a name, a file
extension, a prompt, a title, search lists, and an eight-slot procedure
vector in a fixed order (\code{ref subsystem}):

\begin{pop}
constant myforth_subsystem_procedures =
    {% forth_compile,    ;;; SS_COMPILER   <- the only substantial one
       identfn,          ;;; SS_RESET      (between top-level reads)
       identfn,          ;;; SS_SETUP      (on startup)
       Forth_banner,     ;;; SS_BANNER
       identfn,          ;;; SS_INITCOMP   (compile a standard prelude)
       erase,            ;;; SS_POPARG1    (handle a command-line argument)
       identfn,          ;;; SS_VEDSETUP   (editor; no-op if you skip Ved)
       identfn,          ;;; SS_XSETUP     (X; no-op if you skip graphics)
    %};
\end{pop}

Only \code{SS\_COMPILER} is real work --- it is your tier-2 compiler. The
rest can start as \code{identfn}/\code{erase} stubs. One macro then makes the
language reachable by name from a Pop-11 prompt:

\begin{pop}
define global macro forth;
    "forth" -> sys_compiler_subsystem(`c`);
enddefine;
\end{pop}

Finally, build an image the way the shipped languages are built --- a
\code{mkimage} invocation that loads your sources and installs
\code{mylang.psv} --- and you have \code{poplog mylang}, live subsystem
switching, and your language on every platform Poplog runs on. The four
existing subsystems under \code{pop/lisp/}, \code{pop/pml/}, \code{pop/plog/}
and \code{pop/forth/} are the templates; each is a
\code{*\_subsystem\_procedures} file, a compiler, and a Makefile rule.

\section{The other starting point: extend what is there}

Not every new language needs to be new. The shipped Prolog is 51 files of
ordinary Pop-11, so a specialised variant has an enormous head start: reuse
the term reader, unification and clause store, and replace only the
\emph{control} --- a different search order, tabling, constraint propagation,
a custom resolution rule. Adding built-ins and operators to the existing
engine is hours; swapping the control regime is days to weeks, scaling with
how much of the engine you replace.
